\section{Коэффициенты и ряды Фурье}\label{sec-fourier-coefficients}

\begin{lemma}[Коэффициенты равномерного разложения]\label{lem-orthogonal-coefficients}

Пусть \(\{\varphi_n\}\) --- ортогональная система на \([a,b]\), а ряд \(\sum_{n=1}^\infty c_n\varphi_n(x)\) сходится равномерно на этом отрезке к \(f(x)\). Тогда
\begin{equation}
c_n=\frac{(f,\varphi_n)}{\|\varphi_n\|^2}
=\frac{\displaystyle\int_a^b f(x)\varphi_n(x)\,dx}
{\displaystyle\int_a^b\varphi_n^2(x)\,dx},
\qquad n\ge1.
\label{eq-orthogonal-coefficients}
\end{equation}

\end{lemma}

\begin{proof}

Фиксируем \(k\). Непрерывная функция \(\varphi_k\) ограничена на \([a,b]\). Умножение равномерно сходящегося ряда на неё сохраняет равномерную сходимость. По теореме \ref{thm-series-integral} можно интегрировать почленно:
\[
\begin{aligned}
(f,\varphi_k)
&=\int_a^b\left(\sum_{n=1}^\infty c_n\varphi_n(x)\right)\varphi_k(x)\,dx\\
&=\sum_{n=1}^\infty c_n\int_a^b\varphi_n(x)\varphi_k(x)\,dx
=c_k\|\varphi_k\|^2.
\end{aligned}
\]
При \(n\ne k\) интегралы равны нулю по ортогональности. Делим на положительное число \(\|\varphi_k\|^2\).

\end{proof}

\begin{definition}[Ряд Фурье по ортогональной системе]\label{def-orthogonal-fourier-series}

Пусть \(f\in C([a,b])\). Числа \(c_n\), определённые формулой \eqref{eq-orthogonal-coefficients}, называются коэффициентами Фурье функции \(f\) по системе \(\{\varphi_n\}\). Ряд \(\sum_{n=1}^\infty c_n\varphi_n(x)\) называется её рядом Фурье по этой системе. Записывают
\[
f(x)\sim\sum_{n=1}^\infty c_n\varphi_n(x).
\]
Знак \(\sim\) означает соответствие функции её ряду Фурье; сходимость ряда и равенство его суммы функции требуют отдельного исследования.

\end{definition}

Формула коэффициентов применима и к интегрируемым по Риману функциям \(f\): непрерывность \(\varphi_n\) обеспечивает существование интегралов. Также можно допустить абсолютно сходящиеся несобственные интегралы, если \(f\) локально интегрируема внутри \((a,b)\).

\subsection{Тригонометрический ряд Фурье}\label{sec-fourier-trigonometric}

\begin{definition}[Тригонометрические коэффициенты]\label{def-trigonometric-fourier-series}

Пусть \(f\) абсолютно интегрируема на \([-\pi,\pi]\) в собственном или несобственном смысле. Её тригонометрический ряд Фурье имеет вид
\begin{equation}
f(x)\sim\frac{a_0}{2}+\sum_{n=1}^\infty\bigl(a_n\cos nx+b_n\sin nx\bigr),
\label{eq-trigonometric-fourier-series}
\end{equation}
где
\begin{equation}
\begin{aligned}
a_n&=\frac1\pi\int_{-\pi}^{\pi}f(x)\cos nx\,dx,\qquad n\ge0,\\
b_n&=\frac1\pi\int_{-\pi}^{\pi}f(x)\sin nx\,dx,\qquad n\ge1.
\end{aligned}
\label{eq-trigonometric-coefficients}
\end{equation}

\end{definition}

Значения норм в примере \ref{exm-trigonometric-system} объясняют множитель \(1/\pi\). Коэффициент при постоянной функции равен \((f,1)/\|1\|^2=(2\pi)^{-1}\int_{-\pi}^{\pi}f(x)\,dx=a_0/2\), откуда отдельная запись постоянного члена.

\section{Лемма Римана об осциллирующих интегралах}\label{sec-fourier-riemann}

Обозначение \(f\in R_{\mathrm{loc}}((a,b))\) означает, что \(f\) интегрируема по Риману на каждом замкнутом отрезке, лежащем внутри \((a,b)\). Абсолютная интегрируемость означает конечность \(\int_a^b|f(x)|\,dx\); при особых точках этот интеграл понимается как несобственный.

\begin{lemma}[Римана]\label{lem-riemann-oscillatory}

Пусть \((a,b)\) --- конечный или бесконечный интервал, \(f\in R_{\mathrm{loc}}((a,b))\), а \(\int_a^b|f(x)|\,dx<+\infty\). Тогда
\begin{equation}
\lim_{|\omega|\to\infty}\int_a^b f(x)\sin\omega x\,dx
=\lim_{|\omega|\to\infty}\int_a^b f(x)\cos\omega x\,dx=0.
\label{eq-riemann-oscillatory}
\end{equation}

\end{lemma}

\begin{proof}

\textbf{1. Собственный интеграл на конечном отрезке.}
Пусть \(f\in R([a,b])\) и \(|f(x)|\le C\). Для \(\varepsilon>0\) выбираем разбиение \(a=x_0<x_1<\cdots<x_N=b\), для которого разность верхней и нижней сумм Дарбу удовлетворяет
\[
\sum_{j=1}^N(M_j-m_j)\Delta x_j<\frac{\varepsilon}{2},
\]
где \(M_j=\sup_{[x_{j-1},x_j]}f\), \(m_j=\inf_{[x_{j-1},x_j]}f\), \(\Delta x_j=x_j-x_{j-1}\).
На каждом участке вычитаем постоянную \(m_j\). При \(\omega\ne0\) получаем
\[
\begin{aligned}
\left|\int_a^b f(x)\sin\omega x\,dx\right|
&\le\sum_{j=1}^N\int_{x_{j-1}}^{x_j}|f(x)-m_j|\,dx
+\sum_{j=1}^N|m_j|\left|\int_{x_{j-1}}^{x_j}\sin\omega x\,dx\right|\\
&\le\sum_{j=1}^N(M_j-m_j)\Delta x_j
+\frac{2}{|\omega|}\sum_{j=1}^N|m_j|\\
&<\frac{\varepsilon}{2}+\frac{2NC}{|\omega|}.
\end{aligned}
\]
Разбиение уже фиксировано, поэтому при достаточно большом \(|\omega|\) второе слагаемое меньше \(\varepsilon/2\). Для косинуса доказательство то же: его интеграл по каждому участку также по модулю не превосходит \(2/|\omega|\).

\textbf{2. Абсолютно сходящийся несобственный интеграл.}
Выбираем конечный отрезок \([\alpha,\beta]\subset(a,b)\), чтобы
\[
\int_a^\alpha|f(x)|\,dx+\int_\beta^b|f(x)|\,dx<\frac{\varepsilon}{2}.
\]
На \([\alpha,\beta]\) применяем первую часть. При достаточно большом \(|\omega|\)
\[
\left|\int_a^b f(x)\sin\omega x\,dx\right|
\le\int_a^\alpha|f(x)|\,dx
+\left|\int_\alpha^\beta f(x)\sin\omega x\,dx\right|
+\int_\beta^b|f(x)|\,dx<\varepsilon.
\]
Для косинуса рассуждение аналогично. Если есть конечное число внутренних особых точек, интервал разбиваем в них и применяем доказанное к каждому из полученных интервалов.

\end{proof}

\begin{cor}[Стремление коэффициентов к нулю]\label{cor-fourier-coefficients-zero}

Для абсолютно интегрируемой функции на \([-\pi,\pi]\) коэффициенты \eqref{eq-trigonometric-coefficients} удовлетворяют \(a_n\to0\) и \(b_n\to0\) при \(n\to\infty\).

\end{cor}

Это непосредственное применение леммы \ref{lem-riemann-oscillatory} при \(\omega=n\). Данное свойство само по себе не гарантирует сходимости ряда \eqref{eq-trigonometric-fourier-series}.
